
\chapterwithtoc{Introduction générale}
\label{chap:introduction}

Pendant longtemps, défendre un système d'information a consisté à en garder
la porte. On admet aujourd'hui que l'intrusion finira par se produire, et que la valeur d'une défense se mesure à sa capacité de la détecter, de la comprendre et de la contenir pendant qu'elle a lieu. Au sein d'un centre d'opérations de sécurité, deux équipes sur trois déclarent ne plus suivre le rythme, et les analystes juniors se renouvellent plus vite qu'ils ne se forment~\cite{sans2024}. La recherche a donné un nom à ce phénomène, la \emph{fatigue d'alerte}~\cite{tariq}, et un diagnostic : le
goulot d'étranglement n'est pas la détection, mais la qualification de ce
qui est détecté.

Le référentiel MITRE ATT\&CK~\cite{attack} est devenu la langue commune de cet
exercice ; y cartographier une activité observée reste pourtant un travail
manuel, lent, et tributaire d'une expertise rare. Pour un exploitant
d'infrastructure énergétique, la difficulté redouble : l'incident type
traverse la frontière entre l'informatique de gestion et l'informatique
industrielle (automates, capteurs, protocoles de terrain) deux mondes
qui ne parlent ni les mêmes journaux ni la même matrice ATT\&CK, et dont la
convergence est précisément ce que visent les campagnes les plus avancées.

Les grands modèles de langage (LLM) semblent taillés pour cette tâche. Mais on ne
peut pas encore leur faire confiance : ils \emph{hallucinent}~\cite{ji}, et leur savoir est
en outre figé à la date de leur entraînement, quand le référentiel, lui,
évolue à chaque version. La solution ce n'est pas un modèle plus grand mais un
modèle \emph{ancré}, rattaché au moment de répondre à une source vérifiable et
à jour plutôt qu'à sa seule mémoire. Le \emph{Model Context Protocol} (MCP),
ouvert fin 2024 et confié depuis à une gouvernance
neutre~\cite{mcpintro,aaif}, standardise cet ancrage en donnant au modèle un
accès structuré, déterministe et journalisable à des outils et à des données.

C'est dans ce contexte que s'inscrit le présent essai, qui conçoit un écosystème agentique intégrant le MCP et MITRE ATT\&CK pour
détecter et contextualiser les attaques. Quatre questions le guident :
\begin{enumerate}\setlength{\itemsep}{0pt}\setlength{\parsep}{0pt}\setlength{\topsep}{2pt}
  \item \textbf{Q1 --- Où en sont les modèles sans ancrage ?} Dans quelle mesure
        les modèles de langage actuels savent-ils rattacher au référentiel une
        activité malveillante lue dans des journaux bruts, aux trois niveaux de
        la hiérarchie (tactique, technique, sous-technique) ? Cette capacité se
        maintient-elle lorsque l'on passe de la matrice Enterprise à la matrice
        ICS, c'est-à-dire du journal de système d'information au journal
        d'automate ? La réponse suppose un corpus produit en laboratoire, à
        l'abri de la contamination qui fausse les bancs d'essai publics.
  \item \textbf{Q2 --- Que change l'ancrage déterministe ?} À corpus, consigne
        et conditions d'appel identiques, que gagne un modèle qui interroge un
        référentiel avant de répondre par rapport au même modèle répondant de
        mémoire ? Le gain attendu porte sur trois défauts distincts : les
        identifiants fabriqués, les identifiants périmés, et la confusion entre
        les deux matrices. Encore faut-il les mesurer séparément.
  \item \textbf{Q3 --- Comment construire le serveur qui le permet ?} Quelles
        propriétés un serveur MCP doit-il réunir pour que cet ancrage soit
        exact, reproductible et daté : quelle granularité d'outils, quelle
        version du référentiel, quel traitement des identifiants révoqués,
        quelle trace jointe à chaque réponse ? La littérature fixe des
        conditions d'emploi, elle ne les éprouve pas sur cette tâche.
  \item \textbf{Q4 --- À quel coût en surface d'attaque ?} Ancrer un modèle
        déplace le risque du texte vers l'action : le serveur devient un canal
        d'instruction autant qu'une source de faits. Quelles garanties de
        conception --- lecture seule, moindre privilège, journalisation,
        validation des entrées --- rendent ce déplacement acceptable pour un
        exploitant d'infrastructure essentielle soumis à la LPCE ?
\end{enumerate}